\slajdid{02-s009}
\begin{frame}[label=02-s009]{Način prikazivanja logičkih funkcija i direktna sinteza kombinacionih mreža}
\gusttekst
\begin{columns}[T,onlytextwidth]
\begin{column}{.49\textwidth}
\begin{itemize}\setlength{\itemsep}{0pt}
\item Literal – logička promenljiva ili njena komplementna vrednost (\(A,\bar A\), itd\ldots)
\item Proizvod – jedan literal ili logički proizvod dva ili više literala (\(A,\bar A B,ABC\), itd\ldots)
\item Zbir - jedan literal ili logički zbir dva ili više literala (\(A,\bar A+B,A+B+C\), itd\ldots)
\item Normalni proizvod –proizvod u kojem se promenljiva pojavljuje samo jednom bilo sa svojom pravom, bilo sa svojom komplementnom vrednošću
\item Normalni zbir - zbir u kojem se promenljiva pojavljuje samo jednom bilo sa svojom pravom, bilo sa svojom komplementnom vrednošću
\item Potpuni proizvod – za funkciju sa n promenljivih normalni proizvod u kojem se pojavljuju sve promenljive
\item Potpuni zbir – za funkciju sa n promenljivih normalni zbir u kojem se pojavljuju sve promenljive
\end{itemize}
\end{column}
\begin{column}{.49\textwidth}
\begin{itemize}\setlength{\itemsep}{0pt}
\item Indeks potpunog proizvoda – vrednost n-bitnog binarnog broja dobijenog tako što se u potpunom proizvodu, u kojem su promenljive uređenje na unapred definisan način, svaka promenljiva sa pravom vrednošću zameni binarnom jedinicom a svaka promenljiva sa komplementnom vrednošću binarnom nulom
\item Indeks potpunog zbira – vrednost n-bitnog binarnog broja dobijenog tako što se u potpunom zbiru, u kojem su promenljive uređenje na unapred definisan način, svaka promenljiva sa pravom vrednošću zameni binarnom nulom a svaka promenljiva sa komplementnom vrednošću binarnom jedinicom
\item Zbir proizvoda –logički zbir logičkih proizvoda
\item Proizvod zbirova – logički proizvod logičkih zbirova
\end{itemize}
\end{column}
\end{columns}
\end{frame}
